% ============================================================
%  يُترجَم بـ XeLaTeX  (في Overleaf: Menu → Compiler → XeLaTeX)
%  الخط: Amiri (متوفر في Overleaf، أو حمّله من github.com/aliftype/amiri)
% ============================================================
\documentclass[12pt,a4paper]{article}
\usepackage[a4paper,margin=1.7cm,top=2.3cm,bottom=2.9cm,headheight=15pt,footskip=1.5cm]{geometry}
\usepackage{amsmath,amssymb}
\usepackage{icomma}
\usepackage[english,bidi=default,layout=lists]{babel}
\babelprovide[import,main]{arabic}
\babelfont{rm}[Scale=1.08]{Amiri}
\babelfont{sf}[Scale=1.08]{Amiri}
\usepackage[most]{tcolorbox}
\usepackage{tikz}
\usetikzlibrary{arrows.meta,calc,shadings,patterns}
\usepackage{array}
\usepackage{colortbl}
\usepackage{fancyhdr}
\usepackage{titlesec}
\usepackage{enumitem}
\usepackage{collcell}
\usepackage{needspace}
\usepackage{etoolbox}
\newcolumntype{C}{>{\collectcell\ar}c<{\endcollectcell}}
\pretocmd{\section}{\Needspace{15\baselineskip}}{}{}
\pretocmd{\subsection}{\Needspace{11\baselineskip}}{}{}

\definecolor{cDef}{HTML}{1F4E9C}
\definecolor{cThm}{HTML}{B3261E}
\definecolor{cRem}{HTML}{2E7D32}
\definecolor{cEx}{HTML}{C76B00}
\definecolor{cMeth}{HTML}{6A1B9A}
\definecolor{cSol}{HTML}{00695C}
\definecolor{cSec}{HTML}{1F3A68}

\titleformat{\section}{\Large\bfseries\color{cSec}}{\thesection.}{0.6em}{}[{\color{cSec!60}\titlerule[1pt]}]
\titleformat{\subsection}{\large\bfseries\color{cThm!85!black}}{\thesubsection.}{0.5em}{}
\titlespacing*{\section}{0pt}{14pt}{8pt}
\titlespacing*{\subsection}{0pt}{10pt}{5pt}
\setlength{\parindent}{0pt}
\setlength{\parskip}{4pt}
\setlist{itemsep=2pt,topsep=3pt}

\newtcolorbox{arbox}[3][]{enhanced,
  colback=#2!5,colframe=#2!80!black,arc=2mm,boxrule=0.8pt,
  left=6pt,right=6pt,top=8pt,bottom=5pt,
  fonttitle=\bfseries,coltitle=white,
  title={\foreignlanguage{arabic}{#3}},
  attach boxed title to top right={xshift=-5mm,yshift=-3mm},
  boxed title style={colback=#2!80!black,arc=1.5mm,boxrule=0pt},
  before upper={\selectlanguage{arabic}},
  before skip=8pt,after skip=8pt,breakable,#1}
\newenvironment{BOX}[2]{\begin{otherlanguage}{english}\begin{arbox}{#1}{#2}}{\end{arbox}\end{otherlanguage}}
\newenvironment{tarif}[1][تعريف]{\begin{BOX}{cDef}{#1}}{\end{BOX}}
\newenvironment{khassia}[1][خاصية]{\begin{BOX}{cThm}{#1}}{\end{BOX}}
\newenvironment{molahada}[1][ملاحظة]{\begin{BOX}{cRem}{#1}}{\end{BOX}}
\newenvironment{mithal}[1][مثال]{\begin{BOX}{cEx}{#1}}{\end{BOX}}
\newenvironment{tariqa}[1][طريقة]{\begin{BOX}{cMeth}{#1}}{\end{BOX}}

\newcounter{exo}
\newenvironment{tamrin}[1]{\stepcounter{exo}\begin{BOX}{cSec}{التمرين \theexo: #1}}{\end{BOX}}
\newenvironment{hal}{\Needspace{12\baselineskip}\begin{otherlanguage}{english}\begin{arbox}[breakable]{cSol}{حل التمرين \theexo}}{\end{arbox}\end{otherlanguage}}

\newenvironment{LTR}{\begin{otherlanguage}{english}}{\end{otherlanguage}}
\newcommand{\ar}[1]{\foreignlanguage{arabic}{#1}}

\renewcommand{\dots}{\mathinner{\ldotp\ldotp\ldotp}}
\renewcommand{\ldots}{\mathinner{\ldotp\ldotp\ldotp}}
\renewcommand{\cdots}{\mathinner{\cdotp\cdotp\cdotp}}
\newcommand{\N}{\mathbb{N}}
\newcommand{\Z}{\mathbb{Z}}
\newcommand{\D}{\mathbb{D}}
\newcommand{\Q}{\mathbb{Q}}
\newcommand{\R}{\mathbb{R}}
\newcommand{\pgcd}{\operatorname{PGCD}}
\newcommand{\ppcm}{\operatorname{PPCM}}

% \NLseg{a}{b}{[}{]} ; \NLright{a}{[} ; \NLleft{b}{]}  (interval on a number line)
\newcommand{\NLseg}[4]{\begin{LTR}\begin{tikzpicture}[baseline=-0.6ex]
 \draw[-{Stealth[length=2mm]},gray!70](0,0)--(3.4,0);
 \draw[line width=2pt,cThm!80!black](0.7,0)--(2.7,0);
 \node[cThm!80!black,font=\LARGE] at (0.7,0){$#3$};
 \node[cThm!80!black,font=\LARGE] at (2.7,0){$#4$};
 \node[below=3pt,font=\footnotesize] at (0.7,0){$#1$};
 \node[below=3pt,font=\footnotesize] at (2.7,0){$#2$};
\end{tikzpicture}\end{LTR}}
\newcommand{\NLright}[2]{\begin{LTR}\begin{tikzpicture}[baseline=-0.6ex]
 \draw[-{Stealth[length=2mm]},gray!70](0,0)--(3.4,0);
 \draw[line width=2pt,cThm!80!black](1.3,0)--(3.3,0);
 \node[cThm!80!black,font=\LARGE] at (1.3,0){$#2$};
 \node[below=3pt,font=\footnotesize] at (1.3,0){$#1$};
\end{tikzpicture}\end{LTR}}
\newcommand{\NLleft}[2]{\begin{LTR}\begin{tikzpicture}[baseline=-0.6ex]
 \draw[-{Stealth[length=2mm]},gray!70](0,0)--(3.4,0);
 \draw[line width=2pt,cThm!80!black](0.05,0)--(2.1,0);
 \node[cThm!80!black,font=\LARGE] at (2.1,0){$#2$};
 \node[below=3pt,font=\footnotesize] at (2.1,0){$#1$};
\end{tikzpicture}\end{LTR}}

% ---------- social-media icons (pure TikZ) ----------
\definecolor{cFB}{HTML}{1877F2}
\definecolor{cTG}{HTML}{229ED9}
\definecolor{cIG1}{HTML}{FEDA75}
\definecolor{cIG2}{HTML}{D62976}
\definecolor{cIG3}{HTML}{4F5BD5}
\newcommand{\fbicon}{\begin{LTR}\begin{tikzpicture}[baseline=-0.75ex,x=1em,y=1em]
 \fill[cFB] (0,0) circle (0.55);
 \node[white,inner sep=0,font=\usefont{TU}{lmss}{bx}{n}\small] at (0.04,-0.02) {f};
\end{tikzpicture}\end{LTR}}
\newcommand{\igicon}{\begin{LTR}\begin{tikzpicture}[baseline=-0.75ex,x=1em,y=1em]
 \shade[lower left=cIG1,lower right=cIG2,upper left=cIG2,upper right=cIG3,rounded corners=0.28em] (-0.55,-0.55) rectangle (0.55,0.55);
 \draw[white,line width=0.09em,rounded corners=0.18em] (-0.36,-0.36) rectangle (0.36,0.36);
 \draw[white,line width=0.09em] (0,0) circle (0.17);
 \fill[white] (0.22,0.22) circle (0.05);
\end{tikzpicture}\end{LTR}}
\newcommand{\tgicon}{\begin{LTR}\begin{tikzpicture}[baseline=-0.75ex,x=1em,y=1em]
 \fill[cTG] (0,0) circle (0.55);
 \fill[white] (-0.33,0.0) -- (0.32,0.27) -- (0.2,-0.3) -- (0.0,-0.15) -- (-0.08,-0.28) -- (-0.08,-0.08) -- cycle;
 \fill[cTG!70!black] (-0.08,-0.08) -- (0.25,0.2) -- (0.0,-0.15) -- (-0.08,-0.28) -- cycle;
\end{tikzpicture}\end{LTR}}
% page number: black box, white two-digit number, at the right of the social row
\newcommand{\PageBox}{\begin{LTR}\tikz[baseline=(p.base)]\node[fill=black,text=white,inner sep=2.5pt,minimum width=1.6em,minimum height=1.3em,font=\usefont{TU}{lmss}{bx}{n}\small](p){\ifnum\value{page}<10 0\fi\arabic{page}};\end{LTR}}
% visual left-to-right order: Telegram | Instagram | Facebook | page box
\newcommand{\SocialRow}{\begin{otherlanguage}{english}%
 \ar{\textbf{الأستاذ سوايسية هشام للرياضيات}}\ \tgicon\hfill
 \textbf{prof\_hicham\_math}\ \igicon\hfill
 \ar{\textbf{الأستاذ سوايسية هشام للرياضيات}}\ \fbicon\hspace{2.2em}\PageBox\end{otherlanguage}}

\pagestyle{fancy}
\fancyhf{}
\renewcommand{\headrulewidth}{0.6pt}
\renewcommand{\footrulewidth}{0.4pt}
% \SetHeader{right head: document title}{left head: level -- stream}
\newcommand{\SetHeader}[2]{%
 \fancyhead[C]{\parbox{\textwidth}{\ar{#1\hfill #2}}}%
 \fancyfoot[C]{\parbox{\textwidth}{\small\SocialRow}}}
\usepackage{tkz-tab}
\usepackage{multicol}
% variation / sign table (LTR)
\newenvironment{TAB}{\par\begin{center}\begin{LTR}\begin{tikzpicture}[scale=0.95,every node/.style={font=\small}]}{\end{tikzpicture}\end{LTR}\end{center}}
% grid + axes, clipped: \begin{GR}{xmin}{xmax}{ymin}{ymax}{unit cm}
\newenvironment{GR}[5]{\par\begin{center}\begin{LTR}\begin{tikzpicture}[x=#5cm,y=#5cm,font=\footnotesize]
 \draw[gray!25,very thin] (#1,#3) grid (#2,#4);
 \draw[-{Stealth[length=2mm]}] (#1,0)--(#2+0.3,0) node[below left]{$x$};
 \draw[-{Stealth[length=2mm]}] (0,#3)--(0,#4+0.3) node[below left]{$y$};
 \node[below left] at (0,0){$O$};
 \begin{scope}\clip (#1,#3) rectangle (#2,#4);}
{\end{scope}\end{tikzpicture}\end{LTR}\end{center}}
\newcommand{\Ticks}[2]{\foreach \t in {#1} \node[below,font=\scriptsize] at (\t,0){$\t$}; \foreach \t in {#2} \node[left,font=\scriptsize] at (0,\t){$\t$};}
\newcommand{\bx}[1]{\ensuremath{\colorbox{cDef!12}{$\displaystyle #1$}}}
\newcommand{\TwoCol}[4]{\par\begin{LTR}\noindent\begin{minipage}[t]{#1\textwidth}\selectlanguage{arabic}#2\end{minipage}\hfill\begin{minipage}[t]{#3\textwidth}\selectlanguage{arabic}#4\end{minipage}\end{LTR}\par}
\newcommand{\RL}[4]{\TwoCol{#3}{#4}{#1}{#2}}  % \RL{w1}{right content}{w2}{left content}
\newcommand{\ok}{$\textcolor{cRem}{\checkmark}$}
\newcommand{\limn}{\displaystyle\lim_{n\to+\infty}}
% sub-questions, each on its own line
\newenvironment{sq}{\begin{itemize}[leftmargin=2.4em,itemsep=3pt]}{\end{itemize}}
\newcommand{\qa}{\item[(أ)]}\newcommand{\qb}{\item[(ب)]}\newcommand{\qc}{\item[(جـ)]}\newcommand{\qd}{\item[(د)]}\newcommand{\qe}{\item[(هـ)]}
\newcommand{\pa}{\par\smallskip\textbf{(أ)} }\newcommand{\pb}{\par\smallskip\textbf{(ب)} }\newcommand{\pc}{\par\smallskip\textbf{(جـ)} }\newcommand{\pd}{\par\smallskip\textbf{(د)} }\newcommand{\pe}{\par\smallskip\textbf{(هـ)} }
\newcommand{\QQ}[1]{\par\smallskip\textbf{#1.} }
\newenvironment{tmr}{\stepcounter{exo}\begin{BOX}{cSec}{تمرين \theexo}}{\end{BOX}}
\newcommand{\ABC}[3]{\\[2pt]\begin{tabular}{p{0.28\textwidth}p{0.28\textwidth}p{0.28\textwidth}}(جـ) \ #3 & (ب) \ #2 & (أ) \ #1\end{tabular}}
\newtcbox{\bxtbox}{on line,coltext=cThm,colback=cDef!12,colframe=cDef!12,boxrule=0pt,arc=1pt,boxsep=0pt,left=2pt,right=2pt,top=1pt,bottom=1pt}
\newcommand{\bxt}[1]{\foreignlanguage{english}{\bxtbox{\ar{#1}}}}
\newenvironment{FIG}{\par\begin{center}\begin{LTR}}{\end{LTR}\end{center}}
\renewcommand{\bx}[1]{\ensuremath{\colorbox{cDef!12}{$\displaystyle\color{cThm}#1$}}}
\newcommand{\LW}[1]{{\color{cDef}#1}}
\newcommand{\Car}[1]{\left\{\begin{array}{l}#1\end{array}\right.}
\newcommand{\Lm}[1]{\displaystyle\lim_{x\to#1}}
\newcommand{\LDx}[3]{لدينا: $\Lm{#2}#1=\bx{#3}$ لأنّ: \quad}
\newcommand{\LD}{\par\LDx}
\newcommand{\E}{\par\smallskip}
\relpenalty=10000\binoppenalty=10000\tolerance=9000\emergencystretch=0.3\linewidth
\newenvironment{GRB}[7]{\par\begin{center}\begin{LTR}\begin{tikzpicture}[x=#5cm,y=#6cm,font=\footnotesize]
 \draw[gray!25,very thin,xstep=1,ystep=#7] (#1,#3) grid (#2,#4);
 \draw[-{Stealth[length=2mm]}] (#1,0)--(#2+0.3,0) node[below left]{$x$};
 \draw[-{Stealth[length=2mm]}] (0,#3)--(0,#4+0.3*#7) node[below left]{$y$};
 \node[below left] at (0,0){$O$};
 \begin{scope}\clip (#1,#3) rectangle (#2,#4);}
{\end{scope}\end{tikzpicture}\end{LTR}\end{center}}
\renewenvironment{tamrin}[1]{\stepcounter{exo}\gdef\curt{#1}\begin{BOX}{cSec}{#1}}{\end{BOX}}
\newcommand{\curt}{}
\renewenvironment{hal}{\Needspace{12\baselineskip}\begin{otherlanguage}{english}\begin{arbox}[breakable]{cSol}{حل \curt}}{\end{arbox}\end{otherlanguage}}
\renewcommand{\QQ}[1]{\par\smallskip\textbf{#1.} }
\renewcommand{\ABC}[3]{\\[2pt]\begin{tabular}{p{0.28\textwidth}p{0.28\textwidth}p{0.28\textwidth}}(جـ) \ #3 & (ب) \ #2 & (أ) \ #1\end{tabular}}
\newcommand{\Ss}{$\star\star$}
\SetHeader{\small\bfseries\color{cSec}الدوال الأصلية والتكامل -- تمارين البكالوريا}{\small\color{cSec}السنة الثالثة ثانوي -- تسيير واقتصاد}
\begin{document}
\begin{LTR}
\begin{tcolorbox}[enhanced,colback=cSec!10,colframe=cSec!40,arc=3mm,halign=center,boxrule=0.6pt,top=12pt,bottom=12pt,before upper={\selectlanguage{arabic}}]
{\color{cSec}\fontsize{26}{32}\selectfont\bfseries الدوال الأصلية وحساب التكامل والمساحات}\\[8pt]
{\color{cSec!80!black}\large أسئلة مواضيع البكالوريا (2013 -- 2025) -- بدون حلول (للتدرّب) والتمثيلات البيانية}\\[5pt]
{\color{cSec!80!black}\large السنة الثالثة ثانوي -- شعبة تسيير واقتصاد}
\end{tcolorbox}
\end{LTR}

\begin{BOX}{cMeth}{ما نحتاجه في هذه السلسلة}
\RL{0.48}{\textbf{دالة أصلية:} $F$ دالة أصلية لـ $f$ على $I$ يعني $\LW{F'(x)=f(x)}$ من أجل كل $x$ من $I$.\par
\textbf{التكامل:} $\LW{\displaystyle\int_a^bf(x)\,dx=F(b)-F(a)}$.\par
\textbf{القيمة المتوسطة:} $\LW{\mu=\dfrac1{b-a}\displaystyle\int_a^bf(x)\,dx}$.}{0.48}{\textbf{المساحة:} إذا كان $(C_f)$ فوق $(C_g)$ على $[a;b]$ فإن $\LW{\mathcal A=\displaystyle\int_a^b\big[f(x)-g(x)\big]dx}$ (بوحدة المساحة).\par
\textbf{وحدة المساحة:} إذا كانت الوحدة على المحورين $u$ و $v$ (بالسنتيمتر) فإن $\LW{1\,ua=u\times v\ \text{cm}^2}$.\par
\textbf{قبل كل مساحة:} ندرس الوضع النسبي ، ونرسم الحيز.}
\end{BOX}
\par\smallskip الأسئلة المعلَّمة \Ss\ تحتاج الدالة الأسية أو اللوغاريتمية: تُنجز بعد درسيهما.

% =====================================================================
\section{أسئلة لا تحتاج الدالتين الأسية واللوغاريتمية}

\begin{tamrin}{تمرين من بكالوريا 2022 -- الموضوع الأول}
المستوي منسوب إلى معلم $(O;\vec i,\vec j)$. \ $(C_f)$ و $(C_g)$ التمثيلان البيانيان للدالتين $f$ و $g$ المعرفتين على $\R$ بـ:
\[g(x)=(x+1)^2(x-1)\qquad;\qquad f(x)=ax^2+bx-1\quad\text{\ar{حيث $a$ و $b$ عددان حقيقيان.}}\]
\begin{enumerate}
\item \pa بيّن أنه من أجل كل عدد حقيقي $x$: $g(x)=x^3+x^2-x-1$.
\pb عيّن العددين $a$ و $b$ حتى تكون $g$ دالة أصلية لـ $f$ على $\R$.
\item تحقق أنه من أجل كل عدد حقيقي $x$: $f(x)=(x+1)(3x-1)$.
\item \pa حلّل العبارة $g(x)-f(x)$.
\pb استنتج أن $(C_f)$ و $(C_g)$ يتقاطعان في ثلاث نقط يطلب تعيينها.
\item \textbf{(سؤال إضافي)} ادرس الوضع النسبي لـ $(C_f)$ و $(C_g)$ ، ثم احسب مساحة الحيز المحدد بالمنحنيين والمستقيمين $x=0$ و $x=3$.
\end{enumerate}
\end{tamrin}

\begin{tamrin}{تمرين من بكالوريا 2022 -- الموضوع الأول (اختيار من متعدد)}
عيّن الاقتراح الصحيح الوحيد مع التبرير: \quad قيمة العدد الحقيقي $\displaystyle\int_1^2\Big(1+\frac1{x^2}\Big)dx$ هي:
\ABC{$\dfrac12$}{$\dfrac32$}{$\dfrac23$}
\end{tamrin}

\begin{tamrin}{تمرين من بكالوريا 2022 -- الموضوع الثاني (صحيح أم خطأ)}
$f$ و $g$ الدالتان المعرفتان على $\R$ بـ: $f(x)=(x+1)(3x-3)$ و $g(x)=(x+1)\left(x^2-x-2\right)$.\par
أجب بصحيح أو خطأ مع التعليل: \quad «$g$ هي الدالة الأصلية للدالة $f$ والتي تنعدم عند $-1$».
\end{tamrin}

\begin{tamrin}{تمرين من بكالوريا 2020 -- الموضوع الثاني (اختيار من متعدد)}
$f$ الدالة المعرفة على $\R$ بـ: $f(x)=\dfrac13x^2-2x$. \ عيّن الاقتراح الصحيح الوحيد مع التبرير:
\begin{enumerate}
\item الدالة الأصلية لـ $f$ على $\R$ التي تنعدم من أجل $x=1$ هي الدالة $F$ حيث:
\ABC{$F(x)=x^3-x^2$}{$F(x)=\frac19x^3-x^2$}{$F(x)=\frac19x^3-x^2+\frac89$}
\item القيمة المتوسطة للدالة $f$ على المجال $[0;1]$ هي:
\ABC{$\dfrac19$}{$-\dfrac89$}{$\dfrac89$}
\end{enumerate}
\end{tamrin}

\begin{tamrin}{تمرين من بكالوريا 2021 -- الموضوع الأول (اختيار من متعدد)}
$g$ الدالة المعرفة على $]0;+\infty[$ بـ: $g(x)=2x+1-\dfrac1{x^2}$. \quad دالتها الأصلية $G$ على $]0;+\infty[$ والتي تنعدم من أجل القيمة $1$ معرّفة بـ:
\ABC{$G(x)=x^2+x-\frac1x$}{$G(x)=x^2+x-1-\frac1x$}{$G(x)=x^2+x+\frac1x-3$}
\end{tamrin}

\begin{tamrin}{تمرين من بكالوريا 2021 -- الموضوع الأول (صحيح أم خطأ)}
\RL{0.55}{$g$ الدالة المعرفة على $]0;+\infty[$ بتمثيلها البياني $(C)$ ، و $(T)$ المماس لـ $(C)$ في النقطة $A(1;0)$ (الشكل المقابل).\par
أجب بصحيح أو خطأ مع التبرير:\par
«كل دالة أصلية $G$ للدالة $g$ على $]0;+\infty[$ متزايدة تمامًا على $]0;2]$».}{0.4}{%
\begin{GR}{-0.3}{3.6}{-2.6}{3.4}{0.6}
\Ticks{1,2,3}{-2,-1,1,2,3}
\draw[cDef,very thick,domain=0:3.35,samples=80] plot(\x,{\x*\x*\x-3*\x*\x+2});
\draw[cThm,thick,domain=0:2] plot(\x,{-3*\x+3});
\fill (1,0) circle(2pt) node[above right]{$A$}; \node[cDef] at (3.1,3.1){$(C)$}; \node[cThm] at (1.85,-2.3){$(T)$};
\end{GR}}
\end{tamrin}

\begin{tamrin}{تمرين من بكالوريا 2023 -- الموضوع الأول (اختيار من متعدد)}
القيمة المتوسطة للدالة $x\mapsto3(x+1)^2$ على المجال $[0;1]$ تساوي:
\ABC{$7$}{$14$}{$21$}
\end{tamrin}

\begin{tamrin}{تمرين من بكالوريا 2024 -- الموضوع الثاني (اختيار من متعدد)}
$f$ الدالة المعرفة على $\R$ بـ: $f(x)=8x^3+1$. \quad الدالة الأصلية لـ $f$ على $\R$ والتي تنعدم عند $1$ هي الدالة $F$ حيث:
\ABC{$F(x)=x^4+x-2$}{$F(x)=2x^4+x-3$}{$F(x)=2x^4-2x$}
\end{tamrin}

\begin{tamrin}{تمرين من بكالوريا 2024 -- الموضوع الأول}
\RL{0.55}{المستوي منسوب إلى المعلم المتعامد والمتجانس $(O;\vec i,\vec j)$. \ في الشكل المقابل $(C_f)$ التمثيل البياني لدالة $f$ معرّفة على $\R$ ، والذي يقطع حامل محور الفواصل في نقطتين فاصلتاهما $\alpha$ و $1$ ، و $(T)$ مماس المنحنى في النقطة $A(1;0)$.\par
\textbf{بقراءة بيانية:}
\begin{enumerate}
\item عيّن $f(1)$ و $f'(1)$.
\item $F$ دالة أصلية للدالة $f$ على $\R$. \ حدّد حسب قيم $x$ إشارة $f(x)$ ، ثم استنتج اتجاه تغير الدالة $F$.
\end{enumerate}}{0.4}{%
\begin{GR}{-2}{5}{-3}{2}{0.55}
\Ticks{-1,1,2,3,4}{-2,-1,1}
\draw[cDef,very thick,domain=-1.12:5,samples=120] plot(\x,{(1-\x)*(1-exp(-3*(\x+0.7)))/(1+0.15*(\x-1)*(\x-1))});
\draw[cThm,thick,domain=-1:4] plot(\x,{-\x+1});
\fill (1,0) circle(2pt) node[above right]{$A$}; \fill (-0.7,0) circle(2pt) node[below right]{$\alpha$};
\node[cDef] at (3.3,-0.7){$(C_f)$}; \node[cThm] at (-0.6,1.85){$(T)$};
\end{GR}}
\end{tamrin}

% =====================================================================
\section{دوال أصلية وقيمة متوسطة (\Ss)}

\begin{tamrin}{تمرين من بكالوريا 2013 -- الموضوع الثاني (\Ss\ بعد درس اللوغاريتم)}
\textbf{I)} $g$ الدالة المعرفة على $]0;+\infty[$ بـ: $g(x)=\dfrac{-x^2+x+2}{x^2}$.
\begin{enumerate}
\item عيّن تبعًا لقيم $x$ إشارة $g(x)$.
\item \pa تحقق أنه من أجل كل $x$ من $]0;+\infty[$: $g(x)=-1+\dfrac1x+\dfrac2{x^2}$.
\pb استنتج الدوال الأصلية للدالة $g$ على $]0;+\infty[$.
\end{enumerate}
\textbf{II)} $f$ الدالة المعرفة على $]0;8]$ بـ: $f(x)=3-x-\dfrac2x+\ln x$. \ تحقق أن $f$ هي الدالة الأصلية للدالة $g$ على $]0;8]$ والتي تنعدم عند $1$ ، ثم استنتج اتجاه تغير $f$.
\end{tamrin}

\begin{tamrin}{تمرين من بكالوريا 2014 -- الموضوع الأول (\Ss\ بعد درس اللوغاريتم)}
$f$ الدالة المعرفة على $]0;+\infty[$ بـ: $f(x)=x-1-\dfrac{\ln x}x$. \quad احسب القيمة المتوسطة للدالة $f$ على المجال $[1;3]$.
\end{tamrin}

\begin{tamrin}{تمرين من بكالوريا 2016 -- الموضوع الأول (اختيار من متعدد ، \Ss\ بعد درس الأسية)}
$f$ الدالة المعرفة على $\R$ بـ: $f(x)=e^x+2e^{-x}-3$. \quad القيمة المتوسطة $m$ للدالة $f$ على المجال $[0;2]$ ، مدوّر $m$ إلى الوحدة هو:
\ABC{$1$}{$2$}{$-1$}
\end{tamrin}

\begin{tamrin}{تمرين من بكالوريا 2020 -- الموضوع الأول (\Ss\ بعد درس اللوغاريتم)}
$f$ الدالة المعرفة على $]0;+\infty[$ بـ: $f(x)=\dfrac13x^3-2x+2+x\ln x$ ، و $F$ الدالة المعرفة على $]0;+\infty[$ بـ:
\[F(x)=\frac1{12}x^4-\frac54x^2+2x-8+\frac12x^2\ln x\]
بيّن أن $F$ دالة أصلية للدالة $f$ على $]0;+\infty[$.
\end{tamrin}

\begin{tamrin}{تمرين من بكالوريا 2023 -- الموضوع الأول (اختيار من متعدد ، \Ss\ بعد درس اللوغاريتم)}
$g$ الدالة المعرفة على $]0;+\infty[$ بـ: $g(x)=2x-\dfrac1x$. \quad دالتها الأصلية $G$ على $]0;+\infty[$ التي تنعدم من أجل القيمة $1$ معرّفة بـ:
\ABC{$G(x)=x^2+1-\ln x$}{$G(x)=-x^2+1-\ln x$}{$G(x)=x^2-1-\ln x$}
\end{tamrin}

\begin{tamrin}{تمرين من بكالوريا 2023 -- الموضوع الثاني (اختيار من متعدد ، \Ss\ بعد الدرسين)}
$f$ الدالة المعرفة على $\R$ بـ: $f(x)=\dfrac{2e^x}{e^x+1}$. \quad الدالة الأصلية $F$ لـ $f$ على $\R$ والتي تنعدم من أجل القيمة $0$ معرّفة بـ:
\ABC{$-2\ln\left(e^x+1\right)+\ln4$}{$2\ln\left(e^x+1\right)-\ln4$}{$\ln\Big(\frac{e^x+1}2\Big)$}
\end{tamrin}

\begin{tamrin}{تمرين من بكالوريا 2024 -- الموضوع الثاني (اختيار من متعدد ، \Ss\ بعد الدرسين)}
$g$ الدالة المعرفة على $\R$ بـ: $g(x)=\dfrac{e^x}{1+e^x}$. \quad القيمة المتوسطة للدالة $g$ على المجال $[-\ln2;\ln2]$ هي:
\ABC{$2\ln2$}{$\dfrac1{2\ln2}$}{$\dfrac12$}
\end{tamrin}

\begin{tamrin}{تمرين من بكالوريا 2025 -- الموضوع الثاني (\Ss\ بعد الدرسين)}
\begin{enumerate}
\item $f$ الدالة المعرفة على $\R$ بـ: $f(x)=e^x-x-1$ ، و $F$ الدالة المعرفة على $\R$ بـ: $F(x)=e^x-\dfrac12x^2-x$.\par
احسب $F'(x)$ ، ثم استنتج حساب القيمة المتوسطة للدالة $f$ على $[0;1]$.
\item \pa تحقق أن الدالة $h$ المعرفة على $]0;+\infty[$ بـ: $h(x)=x\ln x-x$ أصلية للدالة $\ln$ على $]0;+\infty[$.
\pb استنتج حساب القيمة المتوسطة للدالة $\ln$ على $[1;e]$.
\end{enumerate}
\end{tamrin}

% =====================================================================
\section{حساب التكامل والمساحات (\Ss)}

\begin{tamrin}{تمرين من بكالوريا 2013 -- الموضوع الأول (\Ss\ بعد الدرسين)}
$f$ الدالة المعرفة على $\R^*$ بـ: $f(x)=2x-1+\dfrac1{e^x-1}$ ، و $(C_f)$ تمثيلها البياني في معلم متعامد ومتجانس. \ نقبل أنه من أجل كل $x\neq0$: $f(x)=2x-2+\dfrac{e^x}{e^x-1}$.\par
احسب العدد $\displaystyle\int_1^2f(x)\,dx$ ، ثم فسّره هندسيًا.
\end{tamrin}

\begin{tamrin}{تمرين من بكالوريا 2015 -- الموضوع الأول (\Ss\ بعد الدرسين)}
$f$ الدالة المعرفة على $\R$ بـ: $f(x)=\dfrac{4e^{-x}}{e^{-x}+1}-3$ ، و $(C_f)$ منحناها في معلم متعامد ومتجانس. \ نقبل أن $f(x)=\dfrac4{e^x+1}-3$ ، وأن $f$ متناقصة تمامًا على $\R$.
\begin{enumerate}
\item جد فاصلة نقطة تقاطع $(C_f)$ مع محور الفواصل.
\item احسب مساحة الحيز المستوي المحدد بالمنحنى $(C_f)$ والمستقيمات التي معادلاتها: $x=0$ ، $x=-\ln3$ و $y=0$.
\end{enumerate}
\end{tamrin}

\begin{tamrin}{تمرين من بكالوريا 2015 -- الموضوع الثاني (\Ss\ بعد درس اللوغاريتم)}
$f$ الدالة المعرفة على $]-1;+\infty[$ بـ: $f(x)=-x+1+3\ln(x+1)$ ، و $(\Gamma)$ تمثيلها البياني. \ نقبل أن $(\Gamma)$ يقطع محور الفواصل في نقطتين فاصلتاهما $\beta\in]-0{,}37;-0{,}36[$ و $\alpha\in]7{,}37;7{,}38[$.
\begin{enumerate}
\item $g$ الدالة المعرفة على $]-1;+\infty[$ بـ: $g(x)=(x+1)\ln(x+1)-x$. \ احسب $g'(x)$ ، ثم استنتج دالة أصلية للدالة $f$.
\item \pa احسب $S$ مساحة الحيز المستوي المحدد بـ $(\Gamma)$ وحامل محور الفواصل والمستقيمين $x=0$ و $x=\alpha$.
\pb تحقق أن $S=\Big(\dfrac12\alpha^2-2\alpha-1\Big)ua$ ، ثم عيّن حصرًا لـ $S$.
\item تنتج ورشة يوميًا $7$ آلاف قطعة على الأكثر. تُنمذج الكلفة الهامشية $C_m$ (الوحدة $1000$ دينار) لإنتاج $x$ ألف قطعة على $[0;7]$ بالدالة $f$: $C_m(x)=f(x)$. \ نرمز بـ $C_T(x)$ إلى الكلفة الإجمالية.
\pa عيّن عبارة $C_T(x)$ علمًا أن الكلفة الإجمالية لإنتاج الألف قطعة الأولى هي $\dfrac52$.
\pb قدّر قيمة الكلفة الإجمالية لإنتاج $7$ آلاف قطعة.
\end{enumerate}
\end{tamrin}

\begin{tamrin}{تمرين من بكالوريا 2016 -- الموضوع الثاني (\Ss\ بعد درس اللوغاريتم)}
$f$ الدالة المعرفة على $]0;+\infty[$ بـ: $f(x)=(2x-4)\ln x$ ، و $(C_f)$ تمثيلها البياني.
\begin{enumerate}
\item عيّن نقط تقاطع $(C_f)$ مع حامل محور الفواصل.
\item $F$ الدالة المعرفة على $]0;+\infty[$ بـ: $F(x)=\left(x^2-4x\right)\ln x-\dfrac12x^2+4x$.
\pa بيّن أن $F$ دالة أصلية للدالة $f$ على $]0;+\infty[$.
\pb احسب مساحة الحيز المستوي المحدد بالمنحنى $(C_f)$ والمستقيمات التي معادلاتها: $y=0$ ، $x=1$ و $x=2$.
\end{enumerate}
\end{tamrin}

\begin{tamrin}{تمرين من بكالوريا 2017 -- الموضوع الأول (\Ss\ بعد درس اللوغاريتم)}
$f$ الدالة المعرفة على $]0;+\infty[$ بـ: $f(x)=x+1-\dfrac{3\ln x}x$ ، و $(C_f)$ تمثيلها البياني ، و $(\Delta)$ المستقيم ذو المعادلة $y=x+1$.
\begin{enumerate}
\item ادرس وضعية $(C_f)$ بالنسبة إلى $(\Delta)$.
\item \pa بيّن أن الدالة $h$ حيث $h(x)=\dfrac12(\ln x)^2$ أصلية للدالة $x\mapsto\dfrac{\ln x}x$ على $]0;+\infty[$.
\pb احسب $S$ مساحة الحيز المستوي المحدد بالمنحنى $(C_f)$ والمستقيمات التي معادلاتها: $x=1$ ، $x=e$ و $y=x+1$.
\end{enumerate}
\end{tamrin}

\begin{tamrin}{تمرين من بكالوريا 2017 (الدورة الاستثنائية) -- الموضوع الأول (\Ss\ بعد درس الأسية)}
$f$ الدالة المعرفة على $\R$ بـ: $f(x)=x^2-x-\dfrac12e^{2x}$ ، و $(C_f)$ تمثيلها البياني في معلم متعامد ومتجانس $(O;\vec i,\vec j)$ حيث $\|\vec i\|=2\ \text{cm}$. \ نقبل أن $f$ متناقصة تمامًا على $\R$ ، وأن المعادلة $f(x)=0$ تقبل حلًا وحيدًا $\alpha$ حيث $-0{,}25<\alpha<-0{,}24$.
\begin{enumerate}
\item احسب بالسنتيمتر المربع المساحة $A(\alpha)$ للحيز المستوي المحدد بالمنحنى $(C_f)$ والمستقيمات التي معادلاتها: $x=0$ ، $x=\alpha$ و $y=0$.
\item تحقق أن: $A(\alpha)=\dfrac13\left(4\alpha^3-12\alpha^2+6\alpha+3\right)\ \text{cm}^2$.
\end{enumerate}
\end{tamrin}

\begin{tamrin}{تمرين من بكالوريا 2018 -- الموضوع الأول (\Ss\ بعد درس اللوغاريتم)}
$f$ الدالة المعرفة على $]-2;8[$ بـ: $f(x)=\ln(x+2)+\ln(-x+8)-\ln16$ ، و $(C_f)$ منحناها في معلم متعامد ومتجانس ، الوحدة البيانية $2\ \text{cm}$.
\begin{enumerate}
\item عيّن نقط تقاطع المنحنى $(C_f)$ مع محوري الإحداثيات.
\item $F$ الدالة المعرفة على $]-2;8[$ بـ: $F(x)=(x+2)\ln(x+2)+(x-8)\ln(-x+8)-2x-x\ln16$. \ بيّن أن $F$ دالة أصلية لـ $f$ على $]-2;8[$.
\item احسب بـ $\text{cm}^2$ مساحة الحيز المستوي المحدد بالمنحنى $(C_f)$ والمستقيمات التي معادلاتها: $y=0$ ، $x=0$ و $x=4$.
\end{enumerate}
\end{tamrin}

\begin{tamrin}{تمرين من بكالوريا 2018 -- الموضوع الثاني (\Ss\ بعد درس الأسية)}
$f$ الدالة المعرفة على $[0;+\infty[$ بـ: $f(x)=x+xe^{-x+1}$ ، و $(C_f)$ تمثيلها البياني ، و $(\Delta)$ المستقيم ذو المعادلة $y=x$. \ نقبل أن $f$ متزايدة تمامًا على $[0;+\infty[$ ، وأن المعادلة $f(x)=4$ تقبل حلًا وحيدًا $\alpha$ حيث $3{,}75<\alpha<3{,}77$.
\begin{enumerate}
\item ادرس وضعية $(C_f)$ بالنسبة إلى $(\Delta)$.
\item $F$ الدالة المعرفة على $[0;+\infty[$ بـ: $F(x)=\dfrac12x^2-(x+1)e^{-x+1}$.
\pa بيّن أن $F$ دالة أصلية للدالة $f$ على $[0;+\infty[$.
\pb أوجد القيمة المضبوطة للعدد $\displaystyle\int_1^4f(x)\,dx$ ، ثم أعط تفسيرًا هندسيًا لهذا العدد.
\item تُنمذج الكلفة الهامشية $C_m$ لإنتاج كمية $q$ (مقدرة بآلاف الوحدات) حيث $0\leqslant q\leqslant7$ بالدالة $f$: $C_m(q)=f(q)$ (مقدرة بملايين الدنانير).
\pa ما هي كمية المنتوج التي من أجلها لا تتجاوز الكلفة الهامشية $4$ ملايين دينار؟
\pb نذكّر أن الكلفة الإجمالية $C_T$ دالة أصلية للكلفة الهامشية. احسب القيمة المتوسطة للكلفة الهامشية عندما تنتج الشركة ما بين $1000$ و $4000$ وحدة.
\end{enumerate}
\end{tamrin}

\begin{tamrin}{تمرين من بكالوريا 2019 -- الموضوع الأول (\Ss\ بعد درس اللوغاريتم)}
$f$ الدالة المعرفة على $\R^*$ بـ: $f(x)=x-\dfrac{x-1}{x^2}$ ، و $(C_f)$ تمثيلها البياني ، و $(\Delta)$ المستقيم ذو المعادلة $y=x$.
\begin{enumerate}
\item ادرس الوضع النسبي للمنحنى $(C_f)$ والمستقيم $(\Delta)$.
\item احسب $A$ مساحة الحيز المستوي المحدد بالمنحنى $(C_f)$ والمستقيمات التي معادلاتها: $y=x$ ، $x=1$ و $x=3$.
\end{enumerate}
\end{tamrin}

\begin{tamrin}{تمرين من بكالوريا 2019 -- الموضوع الثاني (\Ss\ بعد درس الأسية)}
$g$ الدالة المعرفة على $]-\infty;0]$ بـ: $g(x)=2x+6-e^{2x+1}$ ، والمعادلة $g(x)=0$ تقبل حلًا وحيدًا $\alpha$ حيث $-3<\alpha<-2{,}9$.\par
$f$ الدالة المعرفة على $]-\infty;0]$ بـ: $f(x)=-2x^2-12x+e^{2x+1}$ ، و $(C_f)$ تمثيلها البياني في معلم متعامد ، الوحدة $1\ \text{cm}$ على محور الفواصل و $0{,}5\ \text{cm}$ على محور التراتيب. \ نقبل أن $f'(x)=-2g(x)$ ، وأن $f(\alpha)=-2\alpha(\alpha+5)+6$.\par
احسب بدلالة $\alpha$ التكامل $\dfrac12\displaystyle\int_\alpha^0f(x)\,dx$ ، ثم فسّر النتيجة بيانيًا.
\end{tamrin}

\begin{tamrin}{تمرين من بكالوريا 2020 -- الموضوع الثاني (\Ss\ بعد درس اللوغاريتم)}
$f$ الدالة المعرفة على $]0;+\infty[$ بـ: $f(x)=3x-2-\dfrac{2\ln x}{x^2}$ ، و $(C_f)$ تمثيلها البياني في معلم متعامد ومتجانس (وحدة الطول $2\ \text{cm}$) ، و $(\Delta)$ المستقيم ذو المعادلة $y=3x-2$.
\begin{enumerate}
\item ادرس وضعية المنحنى $(C_f)$ بالنسبة إلى المستقيم $(\Delta)$.
\item $H$ الدالة المعرفة على $]0;+\infty[$ بـ: $H(x)=\dfrac{1+\ln x}x$. \ نقبل أن $f$ تقبل قيمة حدية صغرى $f(\alpha)\approx0{,}9$ حيث $0{,}9<\alpha<1$.
\pa بيّن أن $H$ دالة أصلية للدالة $x\mapsto-\dfrac{\ln x}{x^2}$ على $]0;+\infty[$.
\pb احسب بـ $\text{cm}^2$ مساحة الحيز المستوي المحدد بالمنحنى $(C_f)$ وحامل محور الفواصل والمستقيمين $x=1$ و $x=2$.
\end{enumerate}
\end{tamrin}

\begin{tamrin}{تمرين من بكالوريا 2023 -- الموضوع الأول (\Ss\ بعد الدرسين)}
$f$ الدالة المعرفة على $[0;+\infty[$ بـ: $f(x)=x+1-\dfrac3{e^x+1}$ ، و $(C_f)$ تمثيلها البياني (وحدة الطول $2\ \text{cm}$) ، و $(\Delta)$ المستقيم ذو المعادلة $y=x+1$.
\begin{enumerate}
\item ادرس وضعية $(C_f)$ بالنسبة إلى $(\Delta)$.
\item $F$ الدالة المعرفة على $[0;+\infty[$ بـ: $F(x)=3x-3\ln\left(e^x+1\right)$.
\pa تحقق أن $F$ أصلية للدالة $x\mapsto\dfrac3{e^x+1}$ على $[0;+\infty[$.
\pb احسب بالسنتيمتر المربع مساحة الحيز المستوي المحدد بالمنحنى $(C_f)$ والمستقيم $(\Delta)$ والمستقيمين اللذين معادلتاهما $x=0$ و $x=\ln2$.
\end{enumerate}
\end{tamrin}

\begin{tamrin}{تمرين من بكالوريا 2023 -- الموضوع الثاني (\Ss\ بعد درس اللوغاريتم)}
$f$ الدالة المعرفة على $]0;+\infty[$ بـ: $f(x)=x^2-x-\ln x$ ، و $(C_f)$ تمثيلها البياني. \ نقبل أن $f'(x)=\dfrac{(2x+1)(x-1)}x$.
\begin{enumerate}
\item ادرس إشارة $f(x)$ (استعمل جدول تغيرات $f$).
\item $F$ الدالة المعرفة على $]0;+\infty[$ بـ: $F(x)=\dfrac13x^3-\dfrac12x^2+x-x\ln x$.
\pa تحقق أن $F$ أصلية للدالة $f$ على $]0;+\infty[$.
\pb احسب مساحة الحيز المستوي المحدد بالمنحنى $(C_f)$ وحامل محور الفواصل والمستقيمين اللذين معادلتاهما $x=1$ و $x=3$.
\end{enumerate}
\end{tamrin}

\begin{tamrin}{تمرين من بكالوريا 2024 -- الموضوع الأول (\Ss\ بعد درس اللوغاريتم)}
$f$ الدالة المعرفة على $]0;+\infty[$ بـ: $f(x)=1+\dfrac{1+2\ln x}{x^2}$ ، و $(C_f)$ تمثيلها البياني ، و $(\Delta)$ المستقيم ذو المعادلة $y=1$.
\begin{enumerate}
\item ادرس الوضع النسبي للمنحنى $(C_f)$ والمستقيم $(\Delta)$.
\item \pa أثبت أن $H:x\mapsto\dfrac{-3-2\ln x}x$ دالة أصلية للدالة $h:x\mapsto\dfrac{1+2\ln x}{x^2}$ على $]0;+\infty[$.
\pb استنتج $\mathcal A$ مساحة الحيز المستوي المحدد بـ $(C_f)$ و $(\Delta)$ والمستقيمين اللذين معادلتاهما $x=1$ و $x=e$.
\end{enumerate}
\end{tamrin}

\begin{tamrin}{تمرين من بكالوريا 2024 -- الموضوع الثاني (\Ss\ بعد درس الأسية)}
$f$ الدالة المعرفة على $\R$ بـ: $f(x)=x+(x+2)e^{-x}$ ، و $(C_f)$ تمثيلها البياني ، و $(\Delta)$ المستقيم ذو المعادلة $y=x$.
\begin{enumerate}
\item ادرس الوضع النسبي للمنحنى $(C_f)$ والمستقيم $(\Delta)$.
\item \pa بيّن أن الدالة $H:x\mapsto-(x+3)e^{-x}$ هي دالة أصلية للدالة $h:x\mapsto(x+2)e^{-x}$ على $\R$.
\pb احسب $\mathcal A$ مساحة الحيز المستوي المحدد بالمنحنى $(C_f)$ والمستقيمات التي معادلاتها: $y=x$ ، $x=0$ و $x=2$.
\end{enumerate}
\end{tamrin}

\begin{tamrin}{تمرين من بكالوريا 2025 -- الموضوع الأول (\Ss\ بعد درس اللوغاريتم)}
$f$ الدالة المعرفة على $[0;+\infty[$ بـ: $f(0)=0$ ، ومن أجل كل $x$ من $]0;+\infty[$: $f(x)=\dfrac14x^2\left(2\ln x-3\right)$ ، و $(C_f)$ تمثيلها البياني.
\begin{enumerate}
\item ادرس إشارة $f(x)$ على $]0;+\infty[$.
\item $F$ الدالة المعرفة على $]0;+\infty[$ بـ: $F(x)=\dfrac1{36}x^3\left(6\ln x-11\right)$.
\pa بيّن أنه من أجل كل $x$ من $]0;+\infty[$: $F'(x)=f(x)$.
\pb استنتج حساب $\mathcal A$ مساحة الحيز المستوي المحدد بالمنحنى $(C_f)$ والمستقيمات التي معادلاتها: $y=0$ ، $x=1$ و $x=e$.
\end{enumerate}
\end{tamrin}

\begin{tamrin}{تمرين من بكالوريا 2025 -- الموضوع الثاني (\Ss\ بعد درس الأسية)}
$f$ الدالة المعرفة على $\R$ بـ: $f(x)=x+2-xe^x$ ، و $(C_f)$ تمثيلها البياني في معلم متعامد ومتجانس (وحدة الطول $2\ \text{cm}$) ، و $(\Delta)$ المستقيم ذو المعادلة $y=x+2$.
\begin{enumerate}
\item ادرس الوضع النسبي للمنحنى $(C_f)$ والمستقيم $(\Delta)$.
\item $h$ الدالة المعرفة على $\R$ بـ: $h(x)=(1-x)e^x$. \ احسب $h'(x)$ ، ثم استنتج بالسنتيمتر المربع حساب $\mathcal A$ مساحة الحيز المستوي المحدد بالمنحنى $(C_f)$ والمستقيمات التي معادلاتها: $y=x+2$ ، $x=-1$ و $x=0$.
\end{enumerate}
\end{tamrin}

% =====================================================================
\section{الكلفة الإجمالية دالة أصلية للكلفة الهامشية (\Ss)}

\begin{tamrin}{تمرين من بكالوريا 2014 -- الموضوع الثاني (\Ss\ بعد درس الأسية)}
$f$ الدالة المعرفة على $[0;7]$ بـ: $f(x)=6(1-2x)e^{-x}+5$.
\begin{enumerate}
\item \pa عيّن العددين الحقيقيين $a$ و $b$ بحيث تكون الدالة $g$ المعرفة على $[0;7]$ بـ: $g(x)=(ax+b)e^{-x}$ دالة أصلية للدالة $h:x\mapsto6(1-2x)e^{-x}$.
\pb استنتج دالة أصلية للدالة $f$ على $[0;7]$.
\item الكلفة الهامشية $C_m$ لصناعة كمية $x$ (مقدرة بالطن) من منتوج ، حيث $x\in[0;7]$ ، تُنمذج بالدالة $f$: $C_m(x)=f(x)$ (الكلفة مقدرة بملايين الدنانير). نذكّر أن دالة الكلفة الإجمالية دالة أصلية لدالة الكلفة الهامشية.
\pa بيّن أن الكلفة الإجمالية $C_T$ معرفة بـ: $C_T(x)=(12x+6)e^{-x}+5x+k$ حيث $k$ عدد حقيقي.
\pb حدّد قيمة $k$ إذا علمت أن الكلفة الثابتة $2$ مليون دينار (أي $C_T(0)=2$).
\end{enumerate}
\end{tamrin}

\begin{LTR}
\begin{tcolorbox}[enhanced,colback=cSec!10,colframe=cSec!40,arc=3mm,halign=center,boxrule=0.6pt,before upper={\selectlanguage{arabic}}]
{\color{cSec}\Large\bfseries بالتوفيق والنجاح}
\end{tcolorbox}
\end{LTR}
\end{document}
